\documentclass[10pt,a4paper]{amsart}
\usepackage[margin=19mm]{geometry}
\usepackage{amsmath,amssymb,mathtools}
\usepackage[scr=rsfs,cal=euler]{mathalfa}
\usepackage{xfrac}
\usepackage[unicode,hidelinks,bookmarks=false]{hyperref}
\newcommand{\ie}{i.e.\ }
\newcommand{\cf}{cf.\ }
\newcommand{\resp}{respectively\ }
\newcommand{\Subsection}{\subsection}
\providecommand{\nobreakdash}{\nobreak\hbox{-}}
\setlength{\emergencystretch}{2em}
\multlinegap0pt
\usepackage{enumerate}
\usepackage{bm}
\usepackage{amssymb}
\usepackage{tikz}
\usepackage{mathtools}
\newtheorem{cor}{Corollary}
\newtheorem{lem}{Lemma}
\newtheorem{thm}{Theorem}
\DeclareMathOperator{\Gr}{Gr}
\title{\bf Rational Approximations for Reciprocals of Multiple Zeta Values and Trivariate Cauchy Numbers}
\author{{Ce Xu${}^{a,}${} \ \ and Jianqiang Zhao${}^{b,}$}}
\date{Source: arXiv:2609.11072v1 (CC BY 4.0)}
\begin{document}
\maketitle

\noindent\emph{Source and licence.} \bf Rational Approximations for Reciprocals of Multiple Zeta Values and Trivariate Cauchy Numbers. Version arXiv:2609.11072v1 (CC BY 4.0), distributed by arXiv under the Creative Commons Attribution 4.0 International licence, http://creativecommons.org/licenses/by/4.0/, as recorded in the arXiv licence metadata for this exact version and shown on the abstract page https://arxiv.org/abs/2609.11072v1. Full licence notice: \texttt{licenses/arxiv-2609.11072-CC-BY-4.0.txt}.\par\smallskip

\noindent\textit{Editorial scope.} Counted unit: the article's thm thm:GregoryGeneralOrder (source0.tex lines 473--479). The article's complete original proof of it is delivered with it (source0.tex lines 481--627, one \texttt{proof} environment of 147 lines). No mathematics is written, repaired or re-derived: the statement and the proof are the article's own lines, in the article's own notation, and no retained line is substituted or re-flowed.  Retained uncounted as the source-local context the proof needs: lines 308--312; lines 315--317; lines 382--395.  Nothing was omitted from any retained span, so the package carries no omission record.  No retained line contains a citation, so the wrapper carries no bibliography and no bibliography entry.  No retained line contains a figure environment, an image-inclusion command or drawing code, so no image is delivered and the package declares no asset dependency.  Editorial formatting only, in addition: an AMS \texttt{amsart} preamble that restates the article's own declarations and macros used by the retained lines, namely the environments \texttt{cor}, \texttt{lem}, \texttt{thm} on separate counters, together with the standard packages those lines need.  The retained environments print in this document as Corollary 1, Lemma 1, Theorem 1 in order of appearance, whereas the article prints the same items with the numbering of its own full document.  The complete native source of the article as distributed in the arXiv e-print for this exact version is bundled unchanged as \texttt{sources/209985/source0.tex}.
\begin{align} \label{equ:GnEllIntegral}
(-1)^{n+\ell}  G_n^{(\ell)}
=&\,     \int_1^\infty   \left( \sum_{k=1}^{\lceil \ell/2\rceil} (-1)^k \binom{\ell}{2k-1}
               \frac{\log^{\ell-2k+1}(z) \pi^{2k-2}}{(\log^2(z)+\pi^2)^\ell}  \right)\frac{((-1)^{\ell} z^{n-\ell-1}-1)dz}{(z+1)^{n-\ell+1}}
\end{align}

\begin{cor}\label{cor:order2Conj}
We have $G_1^{(2)}=1, G_2^{(2)}=1/12, G_3^{(2)}=0$, and $(-1)^{n} G_n^{(2)}<0$ for all $n\ge 4$.
\end{cor}

\begin{lem}\label{equ:keyLem}
For all positive integer $\ell$, define the polynomial
\begin{equation*}
    H_{\ell}(x)= \sum_{k=1}^{\lceil \ell/2\rceil} (-1)^{k-1} \binom{\ell}{2k-1} x^{\ell-2k+1}.
\end{equation*}
Then we have
\begin{enumerate}
  \item[\upshape{(i)}] $H_{\ell}(x)$ has exactly $\ell-1$ simple real roots. Among these, $\lfloor (\ell-1)/2 \rfloor$ are positive roots.

  \item[\upshape{(ii)}] Let $\rho_\ell$ be the largest root of $H_{\ell}(x)$. Then the sequence $\{\rho_\ell\}_{\ell\ge 3}$ is strictly
  increasing starting with $\rho_3=\sqrt{3}/3$.
\end{enumerate}

\end{lem}

\begin{thm}\label{thm:GregoryGeneralOrder}
For each order $\ell\ge 1$, there is some positive integer $\Gr(\ell)$ such that
\begin{equation*}
(-1)^{n-1} G_n^{(\ell)}>0
\end{equation*}
for all $n\ge \Gr(\ell)$.
\end{thm}

\begin{proof} The case $\ell=1$ is well-known and the case $\ell=2$ is already confirmed by Corollary~\ref{cor:order2Conj} explicitly. We now assume $\ell\ge 3$ and set
\begin{align*}
h(z)=h^{(\ell)}(z):= &\, (-1)^{\ell} \left( \sum_{k=1}^{\lceil \ell/2\rceil} (-1)^k \binom{\ell}{2k-1}
               \log^{\ell-2k+1}(z) \pi^{2k-2}  \right)\cdot  (\log^2(z)+\pi^2)^{-\ell} \\
=&\, (-1)^{\ell-1} \left(\ell \log^{\ell-1}(z)-\binom{\ell}{3} \log^{\ell-3}(z)\pi^2+ \cdots\right) \cdot  (\log^2(z)+\pi^2)^{-\ell}
\end{align*}
and
\begin{equation*}
f_n(z)=f_n^{(\ell)}(z):=\frac{1-(-1)^{\ell} z^{n-\ell-1}}{(z+1)^{n-\ell+1}}.
\end{equation*}
Let $H_\ell(x)$ be the polynomial as defined in Lemma~\ref{equ:keyLem}. Then
\begin{equation*}
h(z)=(-\pi)^{\ell-1}H_{\ell}\big(\log(z)/\pi\big) \cdot (\log^2(z)+\pi^2)^{-\ell}
\end{equation*}
and therefore its largest real root $r=r^{(\ell)}$ satisfies
\begin{equation*}
r=\exp(\pi \rho_\ell)\ge \exp(\pi \rho_3)=\exp(\pi\sqrt{3}/3)\approx 6.13.
\end{equation*}
Hence, $(-1)^{\ell-1} h(z) = h(1/z) >0$ for all $z>r$. By \eqref{equ:GnEllIntegral} we see that for all $n>4$
\begin{equation}\label{equ:2Integrals}
(-1)^{n-1} G_n^{(\ell)}
= \int_1^\infty h(z) f_n(z)\, dz
= \int_1^{r} h(z)  f_n(z)\, dz+\int_0^{1/r} h(z)  f_n(z)\, dz
\end{equation}
by the substitution $z\to 1/z$ in the second integral.

We will proceed to prove that the second integral on the right-hand side of \eqref{equ:2Integrals},
which is always positive, dominates the first. The second integral is easy to estimate
but the first must be cut into two parts and then bounded separately.

We now consider the second integral first. Set
\begin{equation}\label{equ:alpha}
\alpha=\alpha^{(\ell)}:=\int_0^{1/r}  h(z) \, dz>0.
\end{equation}
Note that
\begin{equation*}
     f_n'(z):=\frac{ (-1)^{\ell-1}z^{n-\ell-2}(n-\ell-1-2z)-n+\ell-1}{(z+1)^{n-\ell+2}}.
\end{equation*}
Hence, $f_n'(z)<0$ for all $n>\ell+3$. Indeed, this is obvious if $\ell$ is even. If $\ell$ is odd, then
\begin{equation*}
    z^{n-\ell-2}(n-\ell-1-2z)-n+\ell-1<(n-\ell-1-2z)-n+\ell-1=-2-2z<0.
\end{equation*}
Thus, we get
\begin{equation}\label{equ:2ndIntegralBd}
 \int_0^{1/r} h(z)  f_n(z)  \, dz > f_n(1/r) \int_0^{1/r}  h(z)  \, dz =\alpha f_n(1/r)>0.
\end{equation}

Now we turn to the first integral on the right-hand side of \eqref{equ:2Integrals}.
We have to consider two separate cases depending on the parity of $\ell$.

\medskip\noindent
(i) If $\ell$ is odd then clearly for all $z\in [1, r]$
\begin{equation*}
0< f_n(z)=\frac{z^{n-\ell-1}+1}{(z+1)^{n-\ell+1}}<1
\end{equation*}
and
\begin{equation*}
f_n'(z)=\frac{g_n(z)}{(z+1)^{n-\ell+2}}, \quad\text{where}\quad g_n(z)=z^{n-\ell-2}(n-\ell-1-2z)-n+\ell-1.
\end{equation*}
Since
\begin{equation*}
g_n'(z)=(n-\ell-1)z^{n-\ell-3}(n-\ell-2-2z)>0
\end{equation*}
for all $1\le z\le r< (n-\ell-2)/2$ if $n$ satisfies
\begin{equation}\label{equ:keyRestrict}
n\ge \ell+2r+2,
\end{equation}
$f_n'(z)$ is increasing and has exactly one zero $z_0\in [1,r]$.
We see that $f_n(z)$ attains the minimum at $z_0$ (as $f_n'(1)=-2^{\ell-n}<0$) and $f_n(z)$ achieves maximum at either 1 or $r$.
Moreover, we must have $z_0<2$ since
\begin{equation*}
g_n(2)> 2(n-\ell-5)-n+\ell-1=n-\ell-11>0
\end{equation*}
by \eqref{equ:keyRestrict} as $r>6$. This implies that $f_n(z)$ is increasing on $[2,r]$.
Further, it can be proved easily that
\begin{equation*}
\left(\frac{3}{4}\right)^{n-\ell+1}< \frac14
\end{equation*}
for all $n$ satisfying \eqref{equ:keyRestrict}. Hence,
\begin{equation}\label{equ:fbLargest}
f_n(1)=\frac{1}{2^{n-\ell+1}}<\frac{2^{n-\ell-1}}{3^{n-\ell+1}}<\frac{2^{n-\ell-1}+1}{3^{n-\ell+1}}=f_n(2)<f_n(b)<f_n(r)
\end{equation}
for all $b\in (2,r)$.

We now recall that $h(r)=0$ which implies the existence of some $b\in (2,r)$ such that
\begin{equation*}
\beta_2 (b):= \int_b^r | h(z)| \, dz < \alpha r^2,
\end{equation*}
where $\alpha$ is defined by \eqref{equ:alpha}, which is equivalent to  $\alpha>\beta_2 (b)/r^2$.
Thus,
\begin{align}
\frac {f_n(b)}{f_n(1/r)}=&\, \frac{b^{n-\ell-1}+1}{(b+1)^{n-\ell+1}}\cdot\frac{ (1/r+1)^{n-\ell+1}}{(1/r)^{n-\ell-1}+1}\nonumber\\
=&\, \frac{1}{b^2}  \cdot \frac{(1/b)^{n-\ell-1}+1}{(1/r)^{n-\ell-1}+1} \left(\frac{1/r+1}{1/b+1}\right)^{n-\ell+1}\nonumber\\
<&\,  \frac{2}{b^2}\cdot \left(\frac{1/r+1}{1/b+1}\right)^{n-\ell+1}\nonumber\\
<&\,  \frac{1}{\beta_1 (b)}\left(\alpha-\frac{\beta_2 (b) }{r^2}\right)   \label{equ:geomDecrease}
\end{align}
for all $n\gg 0$ since $1/r+1<1/b+1$, where
\begin{equation*}
\beta_1 (b):= \int_1^b | h(z)| \, dz.
\end{equation*}
By \eqref{equ:fbLargest}, for all sufficiently large $n$, we get
\begin{align*}
\left| \int_1^r h(z) f_n(z)\, dz \right|
< &\,    f_n(b)\beta_1 (b) +  f_n(r)\beta_2 (b) \\
< &\, f_n(1/r) \left(\alpha-\frac{\beta_2 (b)}{r^2} \right)+  f_n(1/r) \frac{\beta_2 (b)}{r^2}
=\alpha f_n(1/r)
\end{align*}
by \eqref{equ:geomDecrease} and the fact that $r^2f_n(r) =f_n(1/r)$.
Combined with \eqref{equ:2Integrals} and \eqref{equ:2ndIntegralBd}, this yields that
\begin{equation}\label{equ:goal}
(-1)^{n-1} G_n^{(\ell)} \ge \alpha f_n(1/r)- \left| \int_1^r h(z) f_n(z)\, dz \right|>0.
\end{equation}


\medskip\noindent
(ii) If $\ell$ is even then  for all $z> 1$
\begin{equation*}
0<F_n(z)=-f_n(z)=\frac{z^{n-\ell-1}-1}{(z+1)^{n-\ell+1}}
<\frac{z^{n-\ell-1}}{(z+1)^{n-\ell+1}}
\le \left(\frac{z}{z+1} \right)^{n-\ell+1}<1.
\end{equation*}
We now show that $F_n(z)$ is increasing on $(1,r]$.
To do this, we notice  that
\begin{equation*}
F_n'(z)=\frac{g_n(z)}{(z+1)^{n-\ell+2}}=\frac{z^{n-\ell-2}(n-\ell-1-2z)+n-\ell+1}{(z+1)^{n-\ell+2}}
\end{equation*}
and
\begin{equation*}
g_n'(z)=(n-\ell-1)z^{n-\ell-3}(n-\ell-2-2z)>0
\end{equation*}
for all $1\le z\le r$ if $n$ satisfies \eqref{equ:keyRestrict}.
Hence, the function $g_n(z)$ is increasing on $[1,r]$. This implies that $g_n(z)>g_n(1)=2n-2\ell-2>0$.
Thus, $F_n'(z)>0$ and therefore $F_n(z)$ is increasing on $[1,r]$ with $F_n(1)=0$.
Hence, by applying the same argument as in  case (i) for odd $\ell$,
after replacing $f_n(z)$ by $F_n(z)$, we may find $2<b<r$ such that $\beta_2 (b)< \alpha r^2$
and
\begin{align*}
\frac {F_n(b)}{f_n(1/r)}=&\, \frac{b^{n-\ell-1}-1}{(b+1)^{n-\ell+1}}\cdot\frac{ (1/r+1)^{n-\ell+1}}{1-(1/r)^{n-\ell-1}}\\
=&\, \frac{1}{b^2}\cdot \frac{1-(1/b)^{n-\ell-1}}{1-(1/r)^{n-\ell-1}} \left(\frac{1/r+1}{1/b+1}\right)^{n-\ell+1}\\
< &\,  \frac{1}{b^2}\cdot \left(\frac{1/r+1}{1/b+1}\right)^{n-\ell+1}\\
<&\, \frac{1}{\beta_1 (b)}\left(\alpha-\frac{\beta_2 (b) }{r^2}\right)
\end{align*}
for all $n\gg 0$ since $1/r+1<1/b+1$. Hence, by exactly the same argument as
in the odd case, we see that \eqref{equ:goal} holds for all sufficiently large $n$.

We have now completed the proof of Theorem~\ref{thm:GregoryGeneralOrder}.
\end{proof}

\end{document}
